% SPDX-License-Identifier: CC-BY-SA-4.0
% Author: Matthieu Perrin
% Part: 
% Section: 
% Sub-section: 
% Frame: 

\begingroup

\begin{frame}[fragile]{L'objet snapshot}

  \begin{description}[Condition de progression :]
  \item[Spécification séquentielle :] Un tableau d'un registre par processus
  \end{description}
  \begin{itemize}
  \item $write(i,v)$ ne peut être appelé que par $p_i$.
  \item On suppose toutes les valeurs écrites différentes
  \end{itemize}
  
  \begin{lstlisting}
    interface Snapshot<T> {
      
    // T values[n] = [0, ..., 0];
      
       T[] snapshot();
    // {
    //    return values;
    // }
    
       void write(int i, T value);
    // {
    //    values[i] = value;
    // }

    }
  \end{lstlisting}

  \begin{description}[Condition de progression :]
  \item[Critère de cohérence :] Linéarisabilité
  \item[Condition de progression :] Wait-freedom
  \end{description}

  \footnoteref{Y. Afek, H. Attiya, D. Dolev, E. Gafni, M. Merritt, N. Shavit. \textit{Atomic snapshots of shared memory.} JACM (1993)}

\end{frame}

\endgroup
\endinput
